% Experimental Setup

\subsection{Datasets}

We evaluate on five tabular benchmarks spanning regression and classification:

\begin{table}[h]
\centering
\small
\begin{tabular}{llrrr}
\toprule
Dataset & Task & $n$ & Features & Target \\
\midrule
Boston Housing & regression & 506 & 13 & MEDV \\
CA Housing & regression & 20{,}640 & 8 & MedHouseVal \\
Concrete & regression & 1{,}030 & 8 & strength \\
Wine (Red) & classification & 1{,}599 & 11 & quality \\
Wine (R+W) & classification & 6{,}497 & 12 & quality \\
\bottomrule
\end{tabular}
\caption{Dataset summary.}
\label{tab:datasets}
\end{table}

Features are standardized to zero mean and unit variance.  Regression targets
are also standardized; predictions are inverse-transformed for evaluation.
The top and bottom 4\% of regression targets are clipped to reduce outlier
sensitivity.  Data is split into 70\% train / 15\% validation / 15\% test
with a fixed random seed for reproducibility.

\subsection{Baselines}

\textbf{Classical ML:} Ridge regression, kernel ridge regression (RBF),
SVR/SVC (RBF), linear/logistic regression, XGBoost, CatBoost.  Boosting
methods use default hyperparameters.

\textbf{Parameter-matched MLP:} A classical MLP whose hidden layer sizes
are chosen so that the total parameter count matches the median quantum
model (${\sim}$200--220 parameters).  This controls for the possibility
that quantum models benefit simply from having more parameters than a
trivially small network.

\subsection{Training Protocol}

All quantum and MLP models are trained with Adam (lr$\,{=}\,0.005$, no
weight decay) for 10{,}000 epochs with full-batch gradient descent.
Best-model checkpointing selects the epoch with lowest validation loss.
Gradient norms are clipped at 1.0 to prevent divergence.

\subsection{Evaluation Metrics}

Regression: $R^2$ (coefficient of determination), RMSE, MAE.
Classification: accuracy, macro-averaged F1 score.
All metrics are reported on the held-out test set using the best-validation
checkpoint.
